\documentclass[../main.tex]{subfiles}
\graphicspath{{\subfix{../img/}}}
\begin{document}
	Gli esami sono generalmente composti da domande a risposta multipla, che però chiedono di motivare la risposta, di seguito saranno elencate le domande, con le relative "crocette", la risposta corretta sarà quella in grassetto, verrà inoltre fornita un eventuale motivazione della risposta.
	\subsection{Domanda 1}
	Una possibile definizione Intelligenza Artificiale generale (Artificial General Intelligence) si colloca in:
	\begin{enumerate}
		\item Sistemi orientati agli obiettivi, progettati per eseguire compiti singoli con abilità che talvolta superano quelle umane. 
		\item \textbf{Sistemi orientati allo sviluppo un modello di intelligenza artificiale, ispirata al funzionamento del cervello umano e pertanto in grado di trovare soluzioni in qualsiasi tipo di situazione, a prescindere dalla natura specifica del problema.}
		\item Sistemi basati su modelli logici/matematici in grado di eseguire compiti il più velocemente possibile. 
	\end{enumerate}
	Motivare adeguatamente la risposta.\\
	\textbf{Risposta}\\
	L'Intelligenza Artificiale Generale (AGI) è un tipo di intelligenza artificiale che ha come obiettivo quello di sviluppare un sistema in grado di risolvere una vasta gamma di compiti simili a quelli che un essere umano può affrontare, senza limitarsi a compiti specifici.\\
	In altre parole, l'AGI dovrebbe essere generalista, capace di adattarsi a qualsiasi tipo di problema, come farebbe una persona, applicando un ragionamento flessibile e una comprensione contestuale che non dipendano dalla specificità di un compito.
	\subsection{Domanda 2}
	Un agente razionale:
	\begin{enumerate}
		\item \textbf{Prende le decisioni sulla base delle conoscenze e dei dati che ha a disposizione al fine di ottenere il miglior risultato atteso tra quelli possibili.}
		\item Esegue sempre la prima azione possibile perché l’importante è essere veloci nell’agire 
		\item Esegue l’azione più giusta secondo i principi della sua coscienza e sulla base dei dati a disposizione. 
	\end{enumerate}
	Motivare adeguatamente la risposta.\\
	\textbf{Risposta}\\
	Un agente razionale è definito in intelligenza artificiale come un'entità che agisce in modo da massimizzare le sue aspettative di successo, cioè le probabilità di ottenere il miglior risultato possibile, dato lo stato attuale delle sue conoscenze e delle risorse disponibili.\\
	L'idea è che l'agente prende decisioni in modo razionale, in funzione delle informazioni che ha, delle sue capacità e dell'ambiente in cui opera.\\
	L'agente non deve necessariamente agire immediatamente, ma deve pianificare la sua azione in modo da ottimizzare il risultato atteso.
	\subsection{Domanda 3}
	Il test di Turing è usato sin dagli inizi dell’Intelligenza artificiale per:
	\begin{enumerate}
		\item \textbf{Valutare sistemi che operano come gli esseri umani}
		\item Valutare sistemi che pensano come gli esseri umani 
		\item Valutare sistemi che agiscono sempre in modo razionale 
	\end{enumerate}
	Motivare adeguatamente la risposta.\\
	\textbf{Risposta}\\
	Il test di Turing è stato pensato per valutare se una macchina è in grado di imitare il comportamento umano in modo tale che un osservatore umano non sia in grado di distinguere se sta interagendo con una macchina o con un altro essere umano.\\
	Il test non si concentra su come la macchina "pensa" o sulla sua coscienza; piuttosto, si concentra sul comportamento osservabile, cioè se la macchina opera in modo simile a un essere umano
	\subsection{Domanda 4}
	Che cosa studia l’explainable AI (XAI)?\\
	\textbf{Risposta}\\
	Studia il motivo per cui è stata presa una decisione dall’AI in modo che i modelli di AI possano essere più interpretabili per gli utenti umani e consentire loro di capire perché il sistema è arrivato a una decisione specifica.
	\subsection{Domanda 5}
	Risoluzione di un problema tramite ricerca:
	\begin{enumerate}
		\item Spiegare come opera un algoritmo di ricerca di una soluzione in uno spazio degli stati, quale è il suo output 
		\item Descrivere quindi i principali algoritmi che adottano strategie non informate. 
	\end{enumerate}
	Queste strategie riescono a trovare soluzioni ottime e complete?\\
	\textbf{Risposta 1}\\
	Lo spazio degli stati è l’insieme di tutti gli stati raggiungibili dallo stato iniziale con una qualunque sequenza di operatori.\\
	E’ caratterizzato da:
	\begin{itemize}
		\item Uno stato iniziale in cui l’agente sa di trovarsi (non noto a priori); 
		\item Un insieme di azioni possibili che sono disponibili da parte dell’agente (Operatori che trasformano uno stato in un altro);
		\item Un cammino è una sequenza di azioni che conduce da uno stato a un altro.
	\end{itemize}
	Il processo di ricerca segue questi passi
	\begin{enumerate}
		\item \textbf{Definizione dello spazio degli stati}: Lo spazio degli stati è una rappresentazione formale di tutte le possibili configurazioni di un problema, a partire da uno stato iniziale e arrivando a uno stato obiettivo.
		\item \textbf{Espansione dei nodi}: L'algoritmo inizia dal nodo iniziale (stato iniziale) e esplora i nodi vicini (stati successivi) mediante azioni che portano a nuovi stati.
		\item \textbf{Strategia di esplorazione}: A seconda dell'algoritmo, la ricerca esplora i nodi in ordine specifico.
		\item \textbf{Condizione di terminazione}: L'algoritmo termina quando raggiunge uno stato che soddisfa la condizione di obiettivo (soluzione) o quando esaurisce lo spazio di ricerca.
		\item \textbf{Output dell'algoritmo}: L'output di un algoritmo di ricerca è solitamente un cammino che collega lo stato iniziale allo stato obiettivo. Se il problema è di tipo di ottimizzazione, l'output potrebbe anche essere la soluzione ottimale.
	\end{enumerate}
	\textbf{Risposta 2}\\
	Non richiedono nessuna conoscenza specifica relativa al problema.
	Non hanno informazioni aggiuntive oltre a quelle già fornite dalla definizione del problema.\\
	Possono solo generare successori e distinguere uno stato obiettivo da uno stato non obiettivo.\\
	Tra i principali algoritmi non informati ci sono:
	\begin{enumerate}
		\item \textbf{Ricerca in ampiezza (Breadth-First Search, BFS}):
		\begin{itemize}
			\item \textbf{Descrizione}: Esplora il grafo degli stati livello per livello, iniziando dal nodo iniziale e visitando tutti i nodi a una distanza di un passo prima di passare ai nodi a distanza maggiore.
			\item \textbf{Proprietà}: È completo (trova una soluzione se esiste), ma non è sempre ottimale (se i costi per passaggio sono diversi, può non trovare la soluzione migliore).
		\end{itemize}
		\item \textbf{Ricerca in profondità (Depth-First Search, DFS)}:
		\begin{itemize}
			\item \textbf{Descrizione}: Esplora profondamente uno dei rami del grafo degli stati, prima di tornare indietro e esplorare gli altri rami.
			\item \textbf{Proprietà}: È completo solo se lo spazio degli stati è finito. Non è garantito che trovi la soluzione ottimale, specialmente in spazi infiniti, e può entrare in loop infiniti se non è implementata correttamente.
		\end{itemize}
		\item \textbf{Ricerca in profondità limitata (Depth-Limited Search, DLS)}:
		\begin{itemize}
			\item \textbf{Descrizione}: È una variante della DFS con un limite di profondità per evitare la ricerca infinita.
			\item \textbf{Proprietà}: È completo solo se esiste una soluzione entro il limite di profondità, ma può non essere ottimale.
		\end{itemize}
		\item \textbf{Ricerca con ritorno su stati precedenti (Iterative Deepening Search, IDS)}:
		\begin{itemize}
			\item \textbf{Descrizione}: Combina la profondità limitata e la ricerca in ampiezza, eseguendo iterazioni della DFS con limiti di profondità crescenti.
			\item \textbf{Proprietà}: È completo e garantisce di trovare la soluzione ottimale (se tutte le azioni hanno lo stesso costo). Tuttavia, ha una complessità computazionale superiore rispetto alla ricerca in ampiezza, poiché riesegue le stesse operazioni più volte.
		\end{itemize}
		\item \textbf{Ricerca Best-First (Best-First Search)}:
		\begin{itemize}
			\item \textbf{Descrizione}: È un algoritmo di ricerca che sceglie quale nodo esplorare in base a un criterio di "priorità", senza sapere se il nodo porterà effettivamente alla soluzione ottimale.
			\item \textbf{Proprietà}: Non è garantito che sia ottimale e potrebbe non essere completo, a meno che non si usino strategie specifiche di gestione dei nodi
		\end{itemize}
	\end{enumerate}
	Le strategie non informate di ricerca sono complete in molti casi (come la BFS e la IDS), ma non sono necessariamente ottimali.\\
	Alcuni di questi algoritmi potrebbero non garantire di trovare la soluzione migliore se non sono applicati in condizioni specifiche (ad esempio, con costi uniformi tra le azioni)
	\subsection{Domanda 6}
	Nella progettazione di un agente logico, con “grounding di un problema” indichiamo:
	\begin{enumerate}
		\item \textbf{La connessione tra i processi di ragionamento logico e l'ambiente reale in cui esiste l'agente}
		\item Il processo di inferenza corretto 
		\item La traduzione della conoscenza iniziale di un problema in un insieme di asserzioni e regole 
		\item La procedura con cui una proposizione A può essere derivata dalla KB base di conoscenza 
	\end{enumerate} 
	Motivare adeguatamente la risposta.\\
	\textbf{Risposta}\\
	Nel contesto della progettazione di un agente logico, il termine "grounding" si riferisce al processo di collegamento delle rappresentazioni simboliche (come proposizioni o concetti) di un agente logico con oggetti concreti o stati del mondo reale. In altre parole, il grounding riguarda la connessione tra il linguaggio formale (come la logica) usato dall'agente e la realtà fisica o l'ambiente in cui l'agente agisce.
	\subsection{Domanda 7}
	Per la costruzione di un agente basato sulla conoscenza si utilizza: 
	\begin{enumerate}
		\item Approccio procedurale 
		\item \textbf{Approccio dichiarativo}
	\end{enumerate}
	Motivare adeguatamente la risposta\\
	\textbf{Risposta}\\
	Un agente basato sulla conoscenza sfrutta l'approccio dichiarativo perché il suo scopo è rappresentare e manipolare la conoscenza in modo esplicito, piuttosto che fornire una sequenza di passi operativi specifici come nell'approccio procedurale.\\ L'approccio dichiarativo permette all'agente di usare la sua base di conoscenza per ragionare e prendere decisioni in modo più flessibile e adattabile.
	\subsection{Domanda 8}
	Tradurre le seguenti frasi inglesi in logica proposizionale: Cercare di conservare il più possibile la struttura della frase:
	\begin{enumerate}
		\item Se e solo se Paolo gioca, Luca va a correre 
		\item La gara sarà finito quando sia Luca che Paolo avranno terminato 10Km di bici e 5 di corsa 
		\item Solo se Luca mette gli occhiali riuscirà a leggere quel cartello 
		\item Usciremo solo se non sarà buio
	\end{enumerate}
	\textbf{Risposta 1}\\
	$P \leftrightarrow L$\\
	\textbf{Risposta 2}\\
	$(L_b \land L_c) \land (P_b \land P_c) \rightarrow G$\\
	\textbf{Risposta 3}\\
	$L_o \rightarrow L_c$\\
	\textbf{Risposta 4}\\
	$U \rightarrow \neg B$
	\subsection{Domanda 9}
	Una possibile definizione si sistemi di Intelligenza Artificiale debole si colloca in:
	\begin{enumerate}
		\item \textbf{Sistemi orientati agli obiettivi, progettati per eseguire compiti singoli con abilità che talvolta superano quelle umane.}
		\item Sistemi orientati allo sviluppo un modello di intelligenza artificiale, ispirata al funzionamento del cervello umano e pertanto in grado di trovare soluzioni in qualsiasi tipo di situazione, a prescindere dalla natura specifica del problema.
		\item Sistemi basati su modelli logici/matematici in grado di eseguire compiti il più velocemente possibile. 
	\end{enumerate}
	Motivare adeguatamente la risposta.\\
	\textbf{Risposta}\\
	L'Intelligenza Artificiale debole (o IA debole) si riferisce a sistemi di intelligenza artificiale progettati per svolgere compiti specifici e ben definiti, come riconoscere oggetti, tradurre lingue o giocare a giochi.\\
	Questi sistemi non sono dotati di consapevolezza o comprensione, ma sono altamente specializzati nel risolvere compiti particolari, spesso con prestazioni che possono superare quelle umane in situazioni limitate e ben definite.
	\subsection{Domanda 10}
	Un agente risolutore di problemi ricerca la soluzione di un problema: 
	\begin{enumerate}
		\item Mentre attua una mossa possibile per le sue capacità dallo stato corrente 
		\item Sulla base del solo stato obiettivo indipendentemente dalla conoscenza delle sue possibilità 
		\item \textbf{Selezionando quella sequenza di azioni possibili per l’agente stesso che porta all’obiettivo dal suo stato corrente}
	\end{enumerate}
	Motivare adeguatamente la risposta\\
	\textbf{Risposta}\\
	Quando un agente affronta un problema, deve identificare una sequenza di azioni che lo porti dallo stato iniziale allo stato obiettivo.\\
	Questo richiede che l'agente esplori lo spazio degli stati possibili, selezionando azioni appropriate ad ogni passo per progredire verso l'obiettivo. Il processo di ricerca di una soluzione implica la valutazione delle possibilità di azione in ogni stato e la determinazione di un percorso che colleghi lo stato iniziale a quello finale.
	\subsection{Domanda 11}
	Un chatbot di AI supera il test di Turing se
	\begin{enumerate}
		\item Risponde a domande su diversi argomenti anche difficili 
		\item \textbf{Imitare le risposte umane in modo ingannevolmente e realistico}
		\item Durante la valutazione adotta un comportamento freddo e calcolatore 
	\end{enumerate}
	Motivare adeguatamente la risposta.\\
	\textbf{Risposta}\\
	La capacità di imitare le risposte umane in modo ingannevole e realistico è fondamentale per superare il test.\\
	Secondo Turing, un chatbot che imita in modo convincente il comportamento umano, a tal punto che il giudice non riesce a distinguere tra una macchina e un umano, può essere considerato in grado di "pensare" o di "simulare" l'intelligenza umana.
	\subsection{Domanda 12}
	Che cosa si intende per “AI black boxes”? \\
	\textbf{Risposta}\\
	I sistemi Machine Learning spesso funzionano come una scatola nera, che fa previsioni e decisioni come fanno gli umani, ma senza essere in grado di comunicare le ragioni per farlo.
	\subsection{Domanda 13}
	Risoluzione di un problema tramite ricerca:
	\begin{enumerate}
		\item Illustrare cosa rappresenta uno spazio degli stati e che cosa la ricerca restituisce come soluzione. 
		\item Spiegare come operano gli algoritmi con strategie informate descrivendo in modo particolare le caratteristiche di una funzione di valutazione
	\end{enumerate}
	\textbf{Risposta 1}\\
	Lo spazio degli stati è l’insieme di tutti gli stati raggiungibili dallo stato iniziale con una qualunque sequenza di operatori. \\
	E’ caratterizzato da:
	\begin{itemize}
		\item Uno stato iniziale in cui l’agente sa di trovarsi (non noto a priori); 
		\item Un insieme di azioni possibili che sono disponibili da parte dell’agente (Operatori che trasformano uno stato in un altro);
		\item Un cammino è una sequenza di azioni che conduce da uno stato a un altro.
	\end{itemize}
	Il processo di ricerca segue questi passi:
	\begin{enumerate}
		\item \textbf{Definizione dello spazio degli stati}: Lo spazio degli stati è una rappresentazione formale di tutte le possibili configurazioni di un problema, a partire da uno stato iniziale e arrivando a uno stato obiettivo.
		\item \textbf{Espansione dei nodi}: L'algoritmo inizia dal nodo iniziale (stato iniziale) e esplora i nodi vicini (stati successivi) mediante azioni che portano a nuovi stati.
		\item \textbf{Strategia di esplorazione}: A seconda dell'algoritmo, la ricerca esplora i nodi in ordine specifico.
		\item \textbf{Condizione di terminazione}: L'algoritmo termina quando raggiunge uno stato che soddisfa la condizione di obiettivo (soluzione) o quando esaurisce lo spazio di ricerca.
		\item \textbf{Output dell'algoritmo}: L'output di un algoritmo di ricerca è solitamente un cammino che collega lo stato iniziale allo stato obiettivo. Se il problema è di tipo di ottimizzazione, l'output potrebbe anche essere la soluzione ottimale.
	\end{enumerate}
	\textbf{Risposta 2}\\
	Gli algoritmi di ricerca con strategie informate (o euristiche) sono utilizzati per risolvere problemi in modo più efficiente rispetto agli algoritmi di ricerca senza informazione (come la ricerca esaustiva, ad esempio, la ricerca a forza bruta o la ricerca in ampiezza).\\
	Questi algoritmi utilizzano informazioni aggiuntive, generalmente sotto forma di una funzione di valutazione, per guidare la ricerca verso soluzioni promettenti.
	\begin{itemize}
		\item \textbf{Euristica}: Fornisce una stima del costo o della distanza rimanente.
		\item \textbf{Direzionalità}: Aiuta a indirizzare la ricerca verso l'obiettivo.
		\item \textbf{Ammissibilità}: Se usata in un algoritmo come A*, deve essere ammissibile per garantire l'ottimalità.
		\item \textbf{Efficienza}: Una buona funzione di valutazione deve essere veloce da calcolare e non comportare un onere computazionale significativo.
	\end{itemize}
	\subsection{Domanda 14}
	Nella progettazione di un agente logico, con “formalizzazione di un problema” indichiamo:
	\begin{enumerate}
		\item 	La connessione tra i processi di ragionamento logico e l'ambiente reale in cui esiste l'agente 
		\item Il processo di inferenza corretto 
		\item \textbf{La traduzione della conoscenza iniziale di un problema in un insieme di asserzioni e regole}
		\item La procedura con cui una proposizione A può essere derivata dalla KB base di conoscenza 
	\end{enumerate}
	Motivare adeguatamente la risposta.\\
	\textbf{Risposta}\\
	La formalizzazione di un problema implica la traduzione della conoscenza (spesso descritta in linguaggio naturale) in una rappresentazione logica formale, che può includere:
	\begin{itemize}
		\item \textbf{Asserzioni} (fatti o dichiarazioni che descrivono lo stato del mondo o le relazioni tra gli oggetti).
		\item \textbf{Regole} (prescrizioni logiche su come la conoscenza può essere manipolata o come nuovi fatti possono essere derivati).
	\end{itemize}
	Questa formalizzazione è fondamentale affinché l'agente possa eseguire il ragionamento logico in modo sistematico e corretto, utilizzando la sua base di conoscenza per risolvere il problema.
	\subsection{Domanda 15}
	Per la costruzione di un agente logico si utilizza:
	\begin{enumerate}
		\item Approccio procedurale
		\item \textbf{Approccio dichiarativo}
	\end{enumerate}
	Motivare adeguatamente la risposta.\\
	\textbf{Risposta}\\
	Questo approccio consente di rappresentare la conoscenza in modo formale, in termini di logica e regole, e di dedurre nuove informazioni tramite inferenza, senza la necessità di specificare dettagliatamente il controllo e le procedure.
	\subsection{Domanda 16}
	Tradurre le seguenti frasi inglesi in logica proposizionale: Cercare di conservare il più possibile la struttura della frase
	\begin{enumerate}
		\item Chiara riesce a lavorare se Marta studia, 
		\item Qualora Chiara si sia in vacanza al mare e Paola sia in vacanza in montagna, la casa in città rimarrà vuota 
		\item Chiara è amica di Marta e viceversa 
		\item Sole se non pioverà, arriveremo in tempo.
	\end{enumerate}
	\textbf{Risposta 1}\\
	$S \rightarrow L$ (Se Marta studia, allora Chiara riesce a lavorare)\\
	\textbf{Risposta 2}\\
	$(M \land P) \rightarrow V$ (Se Chiara è in vacanza al mare e Paola è in vacanza in montagna, allora la casa in città rimarrà vuota)\\
	\textbf{Risposta 3}\\
	$A \land B$ (Chiara è amica di Marta e Marta è amica di Chiara)\\
	\textbf{Risposta 4}\\
	$\neg R \rightarrow T$ (Se non piove, allora arriveremo in tempo)
	\subsection{Domanda 17}
	Tradurre in logica del primo ordine le seguenti affermazioni, commentandone i passi:
	\begin{enumerate}
		\item Le rane sono tutte verdi 
		\item Alcune rane non sono verdi 
		\item Nessuno lettore legge tuti i libri 
		\item In una biblioteca, esiste almeno un libro di storia
	\end{enumerate}
	\textbf{Risposta 1}\\
	$\forall x \left( Rana(x) \to Verde(x) \right)$\\
	\textbf{Risposta 2}\\
	$\exists x \left( Rana(x) \land \neg Verde(x) \right)$\\
	\textbf{Risposta 3}\\
	$\neg \exists x \left( Lettore(x) \land \forall y \left( Libro(y) \to Legge(x, y) \right) \right)$\\
	\textbf{Risposta 4}\\
	$\exists x \left( Libro(x) \land Storia(x) \right)$
	\subsection{Domanda 18}
	In un processo di Data Science si deve estrarre conoscenza dai dati utili per un certo obiettivo. Dato l’obiettivo che si vuole ottenere quali sono i passi che si devono fare per estrarre la conoscenza?\\
	\textbf{Risposta}\\
	I cinque passi per estrarre conoscenza dai dati in un processo di Data Science sono:
	\begin{enumerate}
		\item \textbf{Porre una domanda interessante}: Identificare il problema o l'obiettivo da raggiungere.
		\item \textbf{Ottenere i dati}: Raccogliere o acquisire dati rilevanti per il problema in esame.
		\item \textbf{Esplorare i dati}: Analizzare i dati per capire meglio la loro struttura, rilevanza e qualità.
		\item \textbf{Creare un modello per i dati}: Applicare metodi di modellazione per rappresentare le relazioni nei dati.
		\item \textbf{Comunicare e presentare i risultati}: Condividere i risultati e le conclusioni ottenute in modo chiaro e comprensibile.
	\end{enumerate}
	\subsection{Domanda 19}
	Dopo aver illustrato cosa sono gli ingressi e le uscite di un modello di ML, spiegare la differenza un modello di clustering ed uno di classificazione. Si saprebbe descrivere un algoritmo di classificazione?\\
	\textbf{Risposta}\\
	Ingressi e uscite di un modello di ML:
	\begin{itemize}
		\item Gli ingressi (input) sono i dati grezzi o le caratteristiche rilevanti che vengono forniti al modello per apprendere o fare previsioni. Ad esempio, nel caso della previsione del prezzo di una casa, gli ingressi possono includere la metratura, il numero di stanze, la posizione geografica, ecc.
		\item Le uscite (output) rappresentano i risultati generati dal modello dopo aver elaborato gli ingressi. Possono essere valori continui (per modelli di regressione) o categorie (per modelli di classificazione). Ad esempio, nel caso precedente, l'uscita sarebbe il prezzo previsto della casa.
	\end{itemize}
	Differenza tra modelli di clustering e classificazione:
	\begin{itemize}
		\item Un modello di clustering suddivide i dati in gruppi basati su somiglianze senza conoscere le categorie a priori. \\
		È un approccio non supervisionato.
		\item Un modello di classificazione, invece, assegna etichette predefinite ai dati in base a un processo supervisionato, dove il modello apprende dalle coppie di ingressi e uscite fornite durante l'addestramento.
	\end{itemize}
	Esempio di algoritmo di classificazione:\\
	Un esempio è l'albero di decisione, che utilizza una struttura gerarchica per classificare i dati attraverso una serie di regole logiche; in cui ogni nodo rappresenta una variabile e ogni ramo un possibile valore della variabile, conducendo a una decisione finale (le foglie).
	\subsection{Domanda 20}
	Quando si dice che un modello di ML è andato in overfitting? Ce ne possiamo accorgere durante l’addestramento? \\
	E quando un modello di ML è in underfitting?\\
	\textbf{Risposta}\\
	Un modello è in overfitting quando si adatta troppo ai dati di addestramento, perdendo capacità di generalizzazione ed è rilevabile confrontando prestazioni su addestramento e test.\\
	È in underfitting quando non apprende sufficientemente dai dati, mostrando basse prestazioni su entrambi i set.
	\subsection{Domanda 21}
	Tradurre in logica del primo ordine le seguenti affermazioni, commentandone i passi:
	\begin{enumerate}
		\item Le farfalle sono colorate 
		\item Alcune farfalle non sono rosse. 
		\item Nessuno poliziotto arresta tutti i ladri 
		\item In una pattuglia, esiste almeno un poliziotto graduato. 
	\end{enumerate}
	\textbf{Risposta 1}\\
	$\forall x (Farfalla(x) \to Colorata(x))$\\
	\textbf{Risposta 2}\\
	$\exists x (Farfalla(x) \land \neg Rossa(x))$\\
	\textbf{Risposta 3}\\
	$\neg \exists x (Poliziotto(x) \land \forall y (Ladro(y) \to Arresta(x, y)))$\\
	\textbf{Risposta 4}\\
	$\exists x (Poliziotto(x) \land Graduato(x))$
	\subsection{Domanda 22}
	In un processo di Data Science, la fase della preparazione dei dati richiede in genere molta attenzione e tempo. \\
	Si descriva che cosa si intende per normalizzazione dei dati, selezione delle features e gestione dei dati mancanti.\\
	\textbf{Risposta}\\
	\begin{itemize}
		\item Normalizzazione dei dati: Porta i dati su una scala comune, migliorando la precisione del modello.
		\item Selezione delle features: Determina le variabili più rilevanti per il modello, riducendo la complessità.
		\item Gestione dei dati mancanti: Possono essere trattati imputando valori mancanti (mettere dei valori plausibili al loro posto) o eliminando i dati incompleti.
	\end{itemize}
	\subsection{Domanda 23}
	Dopo aver illustrato cosa sono gli ingressi e le uscite di un modello di ML, spiegare la differenza un modello di regressione ed uno di classificazione.\\
	Si saprebbe descrivere un algoritmo di classificazione?\\
	\textbf{Risposta}\\
	Ingressi e uscite di un modello di ML:
	\begin{itemize}
		\item Gli ingressi (input) sono i dati grezzi o le caratteristiche rilevanti che vengono forniti al modello per apprendere o fare previsioni. Ad esempio, nel caso della previsione del prezzo di una casa, gli ingressi possono includere la metratura, il numero di stanze, la posizione geografica, ecc.
		\item Le uscite (output) rappresentano i risultati generati dal modello dopo aver elaborato gli ingressi. Possono essere valori continui (per modelli di regressione) o categorie (per modelli di classificazione). Ad esempio, nel caso precedente, l'uscita sarebbe il prezzo previsto della casa.
	\end{itemize}
	Differenza tra modelli di regressione e classificazione:
	\begin{itemize}
		\item Un modello di regressione ha come output un valore numerico, che si vuole approssimare tramite una funzione dei dati di input, la variabile di output è continua e può assumere un numero infinito di valori.
		\item Un modello di classificazione, invece, assegna etichette predefinite ai dati in base a un processo supervisionato, dove il modello apprende dalle coppie di ingressi e uscite fornite durante l'addestramento.
	\end{itemize}
	Esempio di algoritmo di classificazione:\\
	Un esempio è l'albero di decisione, che utilizza una struttura gerarchica per classificare i dati attraverso una serie di regole logiche; in cui ogni nodo rappresenta una variabile e ogni ramo un possibile valore della variabile, conducendo a una decisione finale (le foglie).
	\subsection{Domanda 24}
	Nella validazione di un task di classificazione si fa riferimento spesso al parametro accuratezza. Cosa rappresenta e quando può essere una misura ingannevole.\\
	\textbf{Risposta}\\
	L'accuratezza è una metrica che misura la percentuale di previsioni corrette effettuate da un modello di classificazione rispetto al totale delle previsioni e si calcola dividendo il numero di previsioni corrette per il numero totale di previsioni effettuate.\\
	Ad esempio, se un modello classifica correttamente 90 elementi su 100, l'accuratezza sarà del 90\%.\\
	L'accuratezza può risultare ingannevole nei dataset squilibrati, dove una classe è rappresentata in modo molto maggiore rispetto alle altre; ad esempio, se il 90\% dei dati appartiene a una sola classe e il modello predice sempre quella classe, l'accuratezza sarà del 90\%, ma il modello non avrà imparato a distinguere correttamente le classi, in questi casi, metriche come la precisione, il recall e il punteggio F1 sono più utili per valutare le reali prestazioni del modello.
	\subsection{Domanda 25}
	Una possibile definizione Intelligenza Artificiale debole si colloca in: 
	\begin{enumerate}
		\item \textbf{Sistemi che operano in alcuni campi come se fossero intelligenti }
		\item Sistemi che pensano come gli esseri umani 
		\item Sistemi che agiscono razionalmente 
	\end{enumerate}
	Motivare brevemente la risposta.\\
	\textbf{Risposta}\\
	L'intelligenza artificiale debole si riferisce a sistemi che simulano l'intelligenza umana in compiti specifici, senza possedere una vera consapevolezza o coscienza. \\
	Questi sistemi agiscono in modo intelligente solo in determinati contesti, ma non possiedono una comprensione generale del mondo, come potrebbe fare un'intelligenza artificiale forte
	\subsection{Domanda 26}
	Un agente razionale:
	\begin{enumerate}
		\item \textbf{Esegue sempre l’azione corretta in funzione del contesto e delle sue possibilità}
		\item Esegue la prima azione possibile
		\item Esegue l’azione più giusta secondo i suoi principi di etica
	\end{enumerate}
	Motivare brevemente la risposta.\\
	\textbf{Risposta}\\
	Un agente razionale è un'entità che compie azioni che massimizzano il raggiungimento dei suoi obiettivi, tenendo conto dell'ambiente e delle sue capacità.\\
	Non è solo un agente che agisce senza riflettere, ma uno che seleziona la miglior azione in base al contesto.
	\subsection{Domanda 27}
	Problem solving: illustrare cosa fa l’algoritmo di ricerca di una soluzione spiegando che cosa è una strategia di ricerca e i nodi frontiera e come sono relazionati fra di loro\\
	\textbf{Risposta}\\
	Un algoritmo di ricerca di una soluzione esplora lo spazio delle possibili soluzioni partendo dallo stato iniziale fino a raggiungere lo stato finale.\\
	La strategia di ricerca è la tecnica utilizzata per selezionare quale stato esplorare successivamente. \\
	La strategia determina l'ordine in cui i nodi vengono esplorati e può essere basata su costi, profondità, o altri criteri.\\
	I nodi frontiera sono gli stati che sono pronti per essere esplorati (non ancora visitati). \\
	Ogni nodo nella frontiera rappresenta un'opportunità per espandere la ricerca. \\
	Man mano che i nodi vengono esplorati e si trovano soluzioni, i nodi vengono rimossi dalla frontiera e nuovi nodi vengono aggiunti.
	\subsection{Domanda 28}
	Che cosa rappresenta la funzione di valutazione nelle strategie informate?\\
	Fornire un esempio di algoritmo con una strategia informata.\\
	\textbf{Risposta}\\
	La funzione di valutazione in una strategia informata (come la ricerca $A*$) è una funzione che stima quanto un dato stato sia vicino alla soluzione finale, combinando il costo per arrivare a quello stato e una stima del costo rimanente per raggiungere la soluzione.\\
	\underline{Esempio}:\\
	L'algoritmo di ricerca $A*$ utilizza una funzione di valutazione $f(n)=g(n)+h(n)$, dove $g(n)$ è il costo per arrivare al nodo $n$, e $h(n)$ è una stima del costo per arrivare alla soluzione a partire da $n$.
	\subsection{Domanda 29}
	Spiegare che significa, nella progettazione di un agente logico, la seguente affermazione: \\
	Specialized algorithm = specialized knowledge + general‐purpose reasoning\\
	\textbf{Risposta}\\
	L'affermazione implica che un algoritmo specializzato combina due componenti: 
	\begin{itemize}
		\item Conoscenza specializzata, che è il sapere concreto e mirato riguardante un dominio specifico (ad esempio, conoscenze su un particolare campo o problema);
		\item Ragionamento di uso generale, che è un approccio logico per trattare e manipolare tale conoscenza, applicabile in vari contesti
	\end{itemize}
	\subsection{Domanda 30}
	La logica proposizionale è una logica
	\begin{enumerate}
		\item \textbf{Classica} 
		\item Probabilistica 
		\item Temporale 
	\end{enumerate}
	Motivare brevemente la risposta\\
	\textbf{Risposta}\\
	La logica proposizionale è una logica classica in cui le proposizioni sono trattate come entità che possono essere vere o false, senza considerare probabilità, tempo o altre variabili.
	\subsection{Domanda 31}
	Tradurre in logica del primo ordine le seguenti affermazioni, commentandone i passi:
	\begin{enumerate}
		\item Per ogni per ogni fetta di torta, c’è qualcuno che la mangia 
		\item Esiste almeno una torta che non è di colore bianco
	\end{enumerate}
	\textbf{Risposta 1}\\
	\begin{itemize}
		\item \textbf{Traduzione}: $\forall x (FettaDiTorta(x) \to \exists y (Persona(y) \land Mangia(y, x)))$
		\item \textbf{Commento}: Per ogni fetta di torta $x$, esiste qualcuno $y$ (una persona) che mangia quella fetta.
	\end{itemize}
	\textbf{Risposta 2}
	\begin{itemize}
		\item \textbf{Traduzione}: $\exists x (Torta(x) \land \neg ColoreBianco(x))$
		\item \textbf{Commento}: Esiste almeno un oggetto $x$ che è una torta e non ha colore bianco.
	\end{itemize}
	\subsection{Domanda 32}
	Che ruolo ha la conoscenza in un processo di analisi di Business Intelligence\\
	\textbf{Risposta}\\
	La conoscenza è fondamentale in un processo di Business Intelligence (BI), poiché consente di trasformare i dati grezzi in informazioni utili per prendere decisioni strategiche.\\
	La BI si basa su dati storici, analisi predittive e reportistica per generare insight che supportano l'efficacia decisionale e la gestione del business.
	\subsection{Domadna 33}
	Spiegare la differenza fra un modello di ML per il clustering e uno per la classificazione\\
	\textbf{Risposta}
	\begin{itemize}
		\item Clustering: È un tipo di apprendimento non supervisionato dove l'obiettivo è raggruppare dati simili in cluster. Non sono disponibili etichette per il training, quindi l'algoritmo cerca di identificare delle strutture nei dati (ad esempio, K-means).
		\item Classificazione: È un tipo di apprendimento supervisionato in cui l'algoritmo apprende da un set di dati etichettati per poi assegnare etichette a nuovi dati (ad esempio, alberi decisionali, SVM)
	\end{itemize}
	\subsection{Domanda 34}
	Nella matrice di confusione relativo ad un caso di classificazione binaria cosa rappresentano gli elementi definiti come falsi positivi e falsi negativi\\
	\textbf{Risposta}
	\begin{itemize}
		\item Falsi positivi (FP): Sono i casi in cui il modello ha erroneamente classificato una classe negativa come positiva.
		\item Falsi negativi (FN): Sono i casi in cui il modello ha erroneamente classificato una classe positiva come negativa.
	\end{itemize}
	Questi due tipi di errori sono importanti per calcolare metriche come la precisione, il richiamo e l'accuratezza.
	\subsection{Domanda 35}
	Una possibile definizione Intelligenza Artificiale forte si colloca in:
	\begin{enumerate}
		\item Sistemi che imitano gli esseri umani 
		\item \textbf{Sistemi che pensano e sono in grado di prendere decisioni in modo autonomo}
		\item Sistemi che agiscono razionalmente
	\end{enumerate}
	Motivare brevemente la risposta\\
	\textbf{Risposta}\\
	L'Intelligenza Artificiale forte implica che il sistema non solo simuli l'intelligenza umana, ma possieda anche una vera capacità di pensiero e consapevolezza.\\
	Questi sistemi sono in grado di prendere decisioni autonome, come gli esseri umani, comprendendo e ragionando sul mondo che li circonda
	\subsection{Domanda 36}
	Un agente razionale:
	\begin{enumerate}
		\item \textbf{Esegue sempre l’azione corretta in funzione del contesto e delle sue possibilità}
		\item Esegue la prima azione possibile 
		\item Esegue l’azione più giusta secondo i suoi principi di etica 
	\end{enumerate}
	Motivare brevemente la risposta\\
	\textbf{Risposta}\\
	Un agente razionale prende decisioni che massimizzano la probabilità di raggiungere il suo obiettivo, considerando il contesto e le risorse disponibili.\\
	Non si limita a scegliere l'azione prima disponibile, né a seguire principi etici, ma seleziona sempre l'azione che porta al miglior risultato possibile in base alle sue conoscenze e capacità.
	\subsection{Domanda 37}
	Problem solving: illustrare cosa rappresenta uno spazio degli stati e il ruolo degli operatori nella ricerca e costruzione di una soluzione.\\
	\textbf{Risposta}\\
	Lo spazio degli stati è l'insieme di tutte le possibili configurazioni che un sistema può assumere durante il processo di risoluzione di un problema.\\
	Ogni stato rappresenta una possibile situazione intermedia, mentre la soluzione è raggiunta quando si arriva allo stato finale desiderato.\\
	Gli operatori sono le azioni che possono essere applicate per passare da uno stato a un altro. \\
	Ogni operazione modifica lo stato attuale in un nuovo stato, e il loro insieme consente di esplorare lo spazio degli stati alla ricerca della soluzione.\\
	Gli operatori sono quindi fondamentali per costruire la sequenza di azioni necessarie per risolvere un problema.
	\subsection{Domanda 38}
	Parlare brevemente del ruolo dell’euristica nelle strategie informate. \\
	Fornire un esempio di algoritmo con una strategia informata.\\
	\textbf{Risposta}\\
	L'euristica è una funzione che guida la ricerca in modo da esplorare più rapidamente le soluzioni promettenti e ridurre il numero di stati da esaminare. \\
	Nelle strategie informate, l'euristica fornisce una stima di quanto un dato stato sia vicino alla soluzione finale, ottimizzando il processo di ricerca.\\
	\underline{Esempio di algoritmo}: \\
	L'algoritmo di ricerca $A*$ utilizza una funzione euristica per stimare il costo rimanente di un percorso, combinandola con il costo già sostenuto, per determinare la sequenza di azioni migliore da intraprendere per arrivare alla soluzione.
	\subsection{Domanda 39}
	Dovendo progettare un agente logico spiegare brevemente cosa significa tenere separato il motore inferenziale dalla base di conoscenza specifica e che vantaggi si ottengono.\\
	\textbf{Risposta}\\
	Separare il motore inferenziale dalla base di conoscenza permette di creare un agente più flessibile e riutilizzabile.\\
	Il motore inferenziale è responsabile della deduzione delle conclusioni a partire dalle informazioni, mentre la base di conoscenza contiene i fatti e le regole su cui si basa il ragionamento. \\
	Questo approccio consente di aggiornare o sostituire la base di conoscenza senza dover modificare il motore inferenziale, migliorando così la manutenibilità e l'adattabilità dell'agente.
	\subsection{Domanda 40}
	La logica del primo ordine è una logica:
	\begin{enumerate}
		\item Probabilistica 
		\item \textbf{Classica} 
		\item Temporale 
	\end{enumerate}
	Motivare brevemente la risposta\\
	\textbf{Risposta}\\
	La logica del primo ordine è una logica classica perché si basa su principi deterministici e dichiara che ogni affermazione è vera o falsa. \\
	In questa logica, si usano predicati, funzioni e quantificatori per esprimere proposizioni più complesse rispetto alla logica proposizionale
	\subsection{Domanda 41}
	Tradurre in logica del primo ordine le seguenti affermazioni, commentandone i passi:
	\begin{enumerate}
		\item C’è una torta che è mangiata da tutti 
		\item Per ogni fetta di torta c’è un piattino e una forchetta
	\end{enumerate}
	\textbf{Risposta 1}
	\begin{itemize}
		\item Traduzione: $\exists x (\text{Torta}(x) \land \forall y (\text{Persona}(y) \rightarrow \text{Mangia}(y, x)))$
		\item Commento: Esiste un oggetto $x$ che è una torta e per ogni $y$ che è una persona, $y$ mangia $x$.
	\end{itemize}
	\textbf{Risposta 2}
	\begin{itemize}
		\item Traduzione: $\forall x (\text{FettaDiTorta}(x) \rightarrow \exists y \exists z (\text{Piattino}(y) \land \text{Forchetta}(z) \land \text{Usa}(y, x) \land \text{Usa}(z, x)))$
		\item Commento: Per ogni fetta di torta $x$, esistono un piattino $y$ e una forchetta $z$ tali che il piattino $y$ e la forchetta $z$ sono usati per quella fetta.
	\end{itemize}
	\subsection{Domanda 42}
	Su che dati e a che cosa vuole ottenere un processo di Business Intelligence in una azienda\\
	\textbf{Risposta}\\
	Un processo di Business Intelligence raccoglie e analizza dati provenienti da varie fonti aziendali, come transazioni di vendita, dati finanziari, performance dei dipendenti e feedback dei clienti. \\
	L'obiettivo principale è ottenere informazioni utili che supportano la presa di decisioni strategiche, migliorando l'efficienza operativa, individuando tendenze di mercato, e ottimizzando le risorse aziendali.
	\subsection{Domanda 43}
	Spiegare la differenza fra l’uscita di un algoritmo di regressione ed uno di classificazione.\\
	\textbf{Risposta}
	\begin{itemize}
		\item Regressione: L'uscita di un algoritmo di regressione è un valore continuo. \\
		Questo tipo di algoritmo viene utilizzato per prevedere variabili quantitative, come ad esempio il prezzo di una casa o la temperatura di una città.
		\item Classificazione: L'uscita di un algoritmo di classificazione è un valore discreto o un'etichetta che appartiene a una classe definita. \\
		Questo tipo di algoritmo è utilizzato per assegnare un'osservazione a una delle categorie predefinite, come "spam" o "non spam" in un filtro di posta elettronica
	\end{itemize}
	\subsection{Domanda 44}
	Nella validazione di un task di classificazione si fa riferimento spesso al parametro accuratezza.\\
	Cosa rappresenta e quando può essere una misura ingannevole\\
	\textbf{Risposta}\\
	L'accuratezza rappresenta la percentuale di predizioni corrette rispetto al totale delle predizioni effettuate.\\
	È una misura utile quando le classi sono bilanciate.\\
	Tuttavia, può essere ingannevole quando c'è uno squilibrio nelle classi (ad esempio, in un dataset con molte più osservazioni di una classe rispetto all'altra), poiché un modello potrebbe semplicemente prevedere sempre la classe maggioritaria ottenendo comunque un'accuratezza elevata, senza fare previsioni corrette per la classe minoritaria.\\
	In questi casi, altre metriche come precisione, richiamo o F1-score sono più indicative delle performance del modello.
\end{document}